\documentclass[../../../../main.tex]{subfiles}
\begin{document}

座標平面上を次の規則(i)，(ii)に従って1秒ごとに動く点Pを考える。
\begin{enumerate}
  \item最初に，Pは点$(2, \, 1)$にいる。
  \itemある時刻でPが点$(a, \, b)$にいるとき，その1秒後にはPは
\begin{enumerate}
\item[●]確率$\displaystyle\frac{1}{3}$で$x$軸に関して$(a, \, b)$と対称な点
\item[●]確率$\displaystyle\frac{1}{3}$で$y$軸に関して$(a, \, b)$と対称な点
\item[●]確率$\displaystyle\frac{1}{6}$で直線$y=x$に関して$(a, \, b)$と対称な点
\item[●]確率$\displaystyle\frac{1}{6}$で直線$y=-x$に関して$(a, \, b)$と対称な点
\end{enumerate}
にいる。
\end{enumerate}
以下の問いに答えよ。ただし，(1)については，結論のみを書けばよい。
\begin{enumerate}
  \item$P$がとりうる点の座標をすべて求めよ。
  \item$n$を正の整数とする。最初から$n$秒後にPが点$(2, \, 1)$にいる確率と，
最初から$n$秒後にPが点$(-2, \, -1)$にいる確率は等しいことを示せ。
  \item$n$を正の整数とする。最初から$n$秒後にPが点$(2, \, 1)$にいる確率を求めよ。
\end{enumerate}

\end{document}
